\section{Rotation Properties}
\label{sec:rotation}

Cyclic permutation preserves every structural invariant: the same
elements, same size, same sum, and same bounds.

\begin{itemize}
  \item \href{https://github.com/thiagomata/prime-numbers/blob/list-article-v1.0.0/src/main/scala/v1/chapter3/list/properties/RotationProperties.scala}{Same elements}:
    $\operatorname{rotateAt}(L, k).\text{contains}(x) \iff L.\text{contains}(x)$
  \item \href{https://github.com/thiagomata/prime-numbers/blob/list-article-v1.0.0/src/main/scala/v1/chapter3/list/properties/RotationProperties.scala}{Same size and sum}:
    $|\operatorname{rotateAt}(L, k)| = |L|$, $\sum \operatorname{rotateAt}(L, k) = \sum L$
  \item \href{https://github.com/thiagomata/prime-numbers/blob/list-article-v1.0.0/src/main/scala/v1/chapter3/list/properties/RotationProperties.scala}{Bound preservation}:
    $(\forall x \in L,\, x > v) \implies \forall x \in \operatorname{rotateAt}(L, k),\, x > v$,
    and likewise for the upper bound $\forall x \in L,\, x < b$
  \item \href{https://github.com/thiagomata/prime-numbers/blob/list-article-v1.0.0/src/main/scala/v1/chapter3/list/properties/RotationProperties.scala}{Index shift by one}:
    $\operatorname{rotateAt}(L, 1)(i) = L(i + 1)$
\end{itemize}

Rotating a list at index \texttt{k} swaps the front (first \texttt{k}
elements) and back (remaining elements) --- a cyclic permutation.
Rotation preserves every structural invariant: size, sum, element
membership, and bound properties are unchanged, because the multiset of
elements is the same under any cyclic reordering.

Cyclic permutations arise whenever a fixed set of values is viewed from a
different starting position --- for example, shifting a circular buffer
or aligning a periodic sequence. Rotation invariance means no structural
property is lost when the viewing window moves.

\begin{equation*}
\begin{aligned}
\operatorname{rotateAt}(L,\; k) &= \text{back} \mathbin{\texttt{++}} \text{front}
  \quad\text{where } (\text{front},\; \text{back}) = \operatorname{splitAt}(L, k) \\
  &= [L_k, L_{k+1}, \ldots, L_{n-1}, L_0, L_1, \ldots, L_{k-1}]
\end{aligned}
\end{equation*}

\subsection{Permutation Invariants}

Rotation is a permutation of the underlying multiset: the same elements
appear, and every structural quantity derived from the list is
preserved.

\textbf{Membership.} An element belongs to the original list if and only
if it belongs to the rotated list. Since \texttt{append} order doesn't
affect membership, swapping \texttt{front} and \texttt{back} preserves
the element set.

\begin{equation*}
\begin{aligned}
\operatorname{rotateAt}(L, k).\text{contains}(x) &\iff L.\text{contains}(x)
  &&\text{[Q.E.D.]}
\end{aligned}
\end{equation*}

\textbf{Size and sum.} The rotated list has the same length and the same
total sum. Sum over \texttt{append} is additive and commutative.

\begin{equation*}
\begin{aligned}
|\operatorname{rotateAt}(L, k)| &= |L| &&\text{[Same size]} \\
\sum \operatorname{rotateAt}(L, k) &= \sum L &&\text{[Same sum]}
\end{aligned}
\end{equation*}

\textbf{Bound preservation.} If every element of \texttt{L} is strictly
greater than \texttt{v} (or strictly less than \texttt{b}), the same
holds after rotation.

\begin{equation*}
\begin{aligned}
(\forall x \in L,\, x > v) &\implies \forall x \in \operatorname{rotateAt}(L, k),\, x > v \\
(\forall x \in L,\, x < b) &\implies \forall x \in \operatorname{rotateAt}(L, k),\, x < b
\end{aligned}
\end{equation*}

\subsection{Index Shift Under Rotation by One}

Rotating by one position and looking up index $k$ gives the original
list's element at index $k + 1$. This is the lemma that underlies gap
translation in \texttt{ShiftedList}.

\begin{equation*}
\begin{aligned}
\operatorname{rotateAt}(L, 1)(k) = L(k + 1) \quad \text{for } k + 1 < |L| \quad &
\text{[Q.E.D.]}
\end{aligned}
\end{equation*}

Source:
\href{https://github.com/thiagomata/prime-numbers/blob/list-article-v1.0.0/src/main/scala/v1/chapter3/list/properties/RotationProperties.scala}{\texttt{RotationProperties}}

\begin{lstlisting}[style=scala]
def assertRotateContainsForward(
  list: List[BigInt], index: BigInt, x: BigInt
): Boolean = {
  require(index >= 0); require(list.contains(x))
  ListUtils.rotateAt(list, index).contains(x)
}.holds

def assertRotateSameSize(list: List[BigInt], index: BigInt): Boolean = {
  require(index >= 0)
  ListUtils.rotateAt(list, index).size == list.size
}.holds

def assertRotateSameSum(list: List[BigInt], index: BigInt): Boolean = {
  require(index >= 0)
  ListUtils.sum(ListUtils.rotateAt(list, index)) == ListUtils.sum(list)
}.holds

def assertRotateSameLowerBound(
  list: List[BigInt], index: BigInt, value: BigInt
): Boolean = {
  require(index >= 0); require(ListBoundUtils.allGreaterThan(list, value))
  ListBoundUtils.allGreaterThan(ListUtils.rotateAt(list, index), value)
}.holds

def assertRotatedAtIndexPlusOne(list: List[BigInt], k: BigInt): Boolean = {
  require(k + 1 < list.size)
  ListUtils.rotateAt(list, BigInt(1)).apply(k) == list.apply(k + 1)
}.holds
\end{lstlisting}

{\raggedright
These properties are verified in the
\href{https://github.com/thiagomata/prime-numbers/blob/list-article-v1.0.0/src/main/scala/v1/chapter3/list/properties/RotationProperties.scala}{\texttt{RotationProperties}}
module (11 lemmas total). The forward/backward membership pair and the
upper/lower bound pair cover all permutation invariants; the remaining
lemmas (\texttt{assertAppendContainsLeft/\allowbreak Right/\allowbreak
Decompose/\allowbreak Swap}) are structural helpers consumed by the main
rotation proofs.
\par}
